\begingroup
\renewcommand{\arraystretch}{1.45}
\fancyhead[L]{\small ARITHMETISCH-GEOMETRISCHES MITTEL.}
\Needspace{36\baselineskip}
\AGMArticle{15}
[Die Differentialgleichung zweiter Ordnung für die Glieder der Folge des arithmetisch-geometrischen Mittels nimmt eine Form an, in der sie nur Differentiale von Quotienten der Veränderlichen enthält, wenn man mit einer beliebigen Größe $e$ den Ausdruck
$D+\mathrm d(\frac{1}{\Delta}\,\allowbreak \mathrm d\log e)+\frac{1}{\Delta}(\mathrm d\log e)^2$
bildet. Dieser wird nämlich
\Needspace{10\baselineskip}
\[
\begin{aligned}&-\surd(eea^{\mathrm n}b^{\mathrm n})\mathbin{.}\mathrm d
\left(\frac{1}{\Delta a^{\mathrm n}b^{\mathrm n}}\mathbin{.}\mathrm d
\surd\frac{a^{\mathrm n}b^{\mathrm n}}{ee}\right)\\&\quad-\frac{1}{4\Delta}\left(\mathrm d\log\frac{a^{\mathrm n}}{b^{\mathrm n}}\right)^2\end{aligned}
\]
ein Ausdruck, dessen Wert also unabhängig von $n$ ist und sich auch nicht ändert, wenn man $b^{\mathrm n}$ und $c^{\mathrm n}$ oder $c^{\mathrm n}$ und
$a^{\mathrm n}$ statt $a^{\mathrm n}$ und $b^{\mathrm n}$ setzt. Er verwandelt sich für beständig wachsende positive und für negative $n$ in
\Needspace{10\baselineskip}
\[
\begin{gathered}-e\mathrm M(a,b)\mathbin{.}\mathrm d
\left\{\frac{1}{\Delta\mathrm M(a,b)^2}\,\mathrm d\frac{\mathrm M(a,b)}{e}\right\}
\\[7pt]\mbox{und in}\\[7pt]
-e\mathrm M(a,c)\mathbin{.}\mathrm d
\left\{\frac{1}{\Delta\mathrm M(a,c)^2}\,\mathrm d\frac{\mathrm M(a,c)}{e}\right\}\end{gathered}
\]
dagegen für $n$ gleich null und für $a$, $b$, $c$ als besondere Werte von $e$ jeweils in $+\Delta bbcc,\quad-\Delta ccaa,\quad-\Delta aabb.$

\Needspace{24\baselineskip}
\par Setzt man also zur Abkürzung
\Needspace{10\baselineskip}
\[
\begin{gathered}\surd\frac{a}{\mathrm M(a,b)}=p,\\[7pt]
\surd\frac{b}{\mathrm M(a,b)}=q,\\[7pt]
\surd\frac{c}{\mathrm M(a,b)}=r,\\[7pt]
-\pi\frac{\mathrm M(a,b)}{\mathrm M(a,c)}=\log y\end{gathered}
\]
\Needspace{25\baselineskip}
so wird unter Berücksichtigung der Gleichung für $\Delta$ in Artikel 13:
\Needspace{23\baselineskip}
\[
\begin{aligned}
\displaystyle\frac{\Delta}{2}\mathrm M(a,b)^2
&=\displaystyle\frac14\mathrm d\log y\\&=\frac{1}{p^4}\,\mathrm d\log\frac{r}{q}\\&=\frac{1}{q^4}\,\mathrm d\log\frac{r}{p}\\&=\frac{1}{r^4}\,\mathrm d\log\frac{p}{q}\\[8pt]
&=\displaystyle-\frac{pp}{q^4r^4}\,\mathrm d
\left(\frac{1}{\mathrm d\log y}\mathbin{.}\mathrm d\frac{1}{pp}\right)\\&=\frac{qq}{r^4p^4}\,\mathrm d
\left(\frac{1}{\mathrm d\log y}\mathbin{.}\mathrm d\frac{1}{qq}\right)\\&=\frac{rr}{p^4q^4}\,\mathrm d
\left(\frac{1}{\mathrm d\log y}\mathbin{.}\mathrm d\frac{1}{rr}\right)
\end{aligned}
\]
\Needspace{18\baselineskip}
und die Ausdrücke von $-\frac{\Delta}{2}\mathrm M(a,c)^2$ in
\Needspace{10\baselineskip}
\[
\begin{gathered}\surd\frac{a}{\mathrm M(a,c)},\\[7pt]
\surd\frac{c}{\mathrm M(a,c)},\\[7pt]
\surd\frac{b}{\mathrm M(a,c)},\\[7pt]
-\pi\frac{\mathrm M(a,c)}{\mathrm M(a,b)}\end{gathered}
\]
erhalten dieselbe Form.
\Needspace{25\baselineskip}
\par Die Elimination von je zwei der drei Größen $p,q,r$ ergibt
\Needspace{10\baselineskip}
\[
\begin{aligned}&\left\{\frac{1}{p^4}\frac{1}{\mathrm d\log y}\,\mathrm d\log
\left[\frac{16}{p^6}\frac{1}{\mathrm d\log y}\,\mathrm d
\left(\frac{1}{\mathrm d\log y}\mathbin{.}\mathrm d\frac{1}{pp}\right)\right]\right\}^{2}\\&\quad-\frac{16}{p^6}\frac{1}{\mathrm d\log y}\,\mathrm d
\left(\frac{1}{\mathrm d\log y}\mathbin{.}\mathrm d\frac{1}{pp}\right)-1=0\end{aligned}
\]
\SourceLeaf{392}%
als Differentialgleichung sowohl für $p$ als auch für $q$ und für $r$, also für
$\surd{\frac{a}{\mathrm M(a,b)}}$, $\surd{\frac{b}{\mathrm M(a,b)}}$ und
$\surd{\frac{c}{\mathrm M(a,b)}}$; ebenso auch, wenn man darin
$-\pi\frac{\mathrm M(a,c)}{\mathrm M(a,b)}=\log z$ statt $\log y$ gesetzt denkt,
als Differentialgleichung für $\surd{\frac{a}{\mathrm M(a,c)}}$,
$\surd{\frac{c}{\mathrm M(a,c)}}$ und $\surd{\frac{b}{\mathrm M(a,c)}}$.]

\Needspace{28\baselineskip}
\AGMArticle{16}
[Die Darstellung der Quotienten der Größen $a,b,c,\mathrm M(a,b),\mathrm M(a,c)$ durch Reihen, die nach Potenzen von $y^{\frac14}$ oder $z^{\frac14}$ fortschreiten, lässt sich mithilfe der Fundamentalsätze des hier zu untersuchenden Algorithmus zum Beispiel auf folgende Weise ausführen.

Nach der Definition von $y$ und dem in Artikel 12 für den rückwärts verlängerten Algorithmus aufgestellten Satz gilt:
\Needspace{10\baselineskip}
\[
-\pi\frac{\mathrm M(a,b)}{\mathrm M(a,c)}
=-\pi\frac{\mathrm M(a^{\mathrm n},b^{\mathrm n})}{2^n\mathrm M(a^{\mathrm n},c^{\mathrm n})}=\log y
\]
Also ist nach den Gleichungen in Artikel 12 der Grenzwert von\hfill\\{}
$\frac12y^{-2^{n-2}}\cdot\surd{\frac{c^{\mathrm n}}{\mathrm M(a^{\mathrm n},b^{\mathrm n})}}$ für ein beständig wachsendes $n$ oder, was dasselbe ist, der Grenzwert von $\frac12y^{-\frac14}r(y)$,
wobei $r(y)$ statt $\surd{\frac{c}{\mathrm M(a,b)}}$ gesetzt ist, für positives $y^{\frac14}$, das bis auf null abnimmt, gleich eins.

\Needspace{27\baselineskip}
Bezeichnen wir ferner $\surd{\frac{a}{\mathrm M(a,b)}}$ und $\surd{\frac{b}{\mathrm M(a,b)}}$ durch $p(y)$ und $q(y)$, so folgt aus
$aa=bb+cc$ und den beiden Gleichungen für $\surd{\mathrm M(a,b)}$ in Artikel 12
\Needspace{10\baselineskip}
\[
\begin{aligned}p(y)&=1+r(y^4)+r(y^{16})+r(y^{64})\\&\quad+.\;.+r(y^{2^{2n}})+.\;.\end{aligned}
\]
\Needspace{10\baselineskip}
\[
\begin{aligned}q(y)&=1-r(y^4)+r(y^{16})+r(y^{64})\\&\quad+.\;.+r(y^{2^{2n}})+.\;.\end{aligned}
\]
\Needspace{10\baselineskip}
\[
\begin{aligned}r(y)^4&=p(y)^4-q(y)^4\end{aligned}
\]
\Needspace{25\baselineskip}
Die Reihen für $p,q,r$, die nach ganzen Potenzen von $y^{\frac14}$ fortschreiten und diese Bedingungen erfüllen, findet man, soweit man die Entwicklung ausführt, von der Form:
\Needspace{10\baselineskip}
\[
\begin{aligned}p(y)&=1+2y+2y^4+2y^9\\&\quad+.\;.+2y^{nn}+.\;.\;.\end{aligned}
\]
\Needspace{10\baselineskip}
\[
\begin{aligned}q(y)&=1-2y+2y^4-2y^9\\&\quad+.\;.\pm2y^{nn}\mp.\;.\;.\end{aligned}
\]
\Needspace{10\baselineskip}
\[
\begin{aligned}r(y)&=2y^{\frac14}+2y^{\frac94}+2y^{\frac{25}4}+2y^{\frac{49}4}\\&\quad+.\;.+2y^{(n+\frac12)^2}+.\;.\;.\end{aligned}
\]

\Needspace{29\baselineskip}
\SourceLeaf{393}%
Dass das hier angedeutete Gesetz zur Bildung der Glieder allgemein gültig ist, scheint \textsc{Gauss} unter anderem auch auf folgende Weise bewiesen zu haben. Neben den entsprechenden Gleichungen, die sich auf Reihen mit zwei Argumenten beziehen und weiter unten in den Artikeln 23 und 25 ihren Platz finden werden, hat \textsc{Gauss} sich die folgende Aufzeichnung gemacht:]

\begingroup\itshape
Zur Theorie der Zerlegung der Zahlen in vier Quadrate.

Der Satz: Das Produkt zweier Summen von vier Quadraten ist selbst eine Summe von vier Quadraten, wird am einfachsten so dargestellt:

Es seien $l,m,\lambda,\mu,\lambda',\mu'$ sechs komplexe Zahlen, sodass $\lambda,\lambda'$ und $\mu,\mu'$ jeweils konjugiert sind.
Mit $\mathrm N$ bezeichne man die Norm. Dann ist
\Needspace{10\baselineskip}
\[
(\mathrm Nl+\mathrm Nm)(\mathrm N\lambda+\mathrm N\mu)
=\mathrm N(l\lambda+m\mu)+\mathrm N(l\mu'-m\lambda')
\]
\par\endgroup
\Needspace{20\baselineskip}
[und somit auch
\Needspace{18\baselineskip}
\[
\begin{aligned}
&\{\mathrm N(n+in_{\prime})+\mathrm N(n_{\prime\prime}+in_{\prime\prime\prime})\}
\\&\quad\{\mathrm N(1-i)+\mathrm N(1+i)\}\\
&\qquad=\mathrm N\{(n+n_{\prime}+n_{\prime\prime}-n_{\prime\prime\prime})
\\&\qquad\quad+i(-n+n_{\prime}+n_{\prime\prime}+n_{\prime\prime\prime})\}\\
&\qquad\quad+\mathrm N\{(n+n_{\prime}-n_{\prime\prime}+n_{\prime\prime\prime})
\\&\qquad\quad-i(+n-n_{\prime}+n_{\prime\prime}+n_{\prime\prime\prime})\}
\end{aligned}
\]
Hieraus lassen sich leicht die beiden folgenden Sätze ableiten. In ihnen beziehen sich verschiedene Darstellungen einer Zahl als Summe von vier Quadratzahlen auf die verschiedenen Wertesysteme der vier Wurzeln unter Berücksichtigung sowohl der Vorzeichen als auch der Reihenfolge der Wurzeln; ferner wird die Null zu den geraden Zahlen gerechnet.

\Needspace{10\baselineskip}
Ist das Vierfache einer Zahl von der Form $4k+1$ durch vier ungerade Quadratzahlen darstellbar, so ist die Zahl selbst halb so oft durch eine ungerade und drei gerade Quadratzahlen darstellbar. Umgekehrt ist, wenn eine Zahl auf die letztere Weise darstellbar ist, ihr Vierfaches doppelt so oft auf die erstere Weise darstellbar.

\Needspace{10\baselineskip}
Ist das Vierfache einer Zahl von der Form $4k+3$ durch vier ungerade Quadratzahlen darstellbar, so ist die Zahl selbst halb so oft durch eine gerade und drei ungerade Quadratzahlen darstellbar. Umgekehrt ist, wenn eine Zahl auf die letztere Weise darstellbar ist, ihr Vierfaches doppelt so oft auf die erstere Weise darstellbar.

\Needspace{21\baselineskip}
Mithilfe dieser beiden Sätze lässt sich unmittelbar beweisen, dass die obigen Reihen der Gleichung $r(y)^4=p(y)^4-q(y)^4$ genügen, und ferner durch \SourceLeaf{394}%
die wiederholte Anwendung dieser Beziehung, dass, wenn man nach dem Schema
\Needspace{10\baselineskip}
\[
p+q=2p'',\quad p-q=2r'',\quad(p'')^4-(r'')^4=(q'')^4
\]
aus den Werten der Quadrate der beiden ersten Reihen, nämlich $pp$ und $qq$, als Anfangsgliedern die Glieder des Algorithmus eines arithmetisch-geometrischen Mittels für positive gerade Indizes bildet, auch
\Needspace{10\baselineskip}
\[
p^{2\mathrm n}=p(y^{2^{2n}}),\quad q^{2\mathrm n}=q(y^{2^{2n}}),\quad r^{2\mathrm n}=r(y^{2^{2n}})
\]
gilt. Durch Übergang zum Grenzwert von $n$ ergibt sich also nach Artikel 12:
\Needspace{10\baselineskip}
\[
\mathrm M(pp,qq)=1,\qquad
\frac\pi2\frac{\mathrm M(pp,qq)}{\mathrm M(pp,rr)}=-\tfrac12\log y\ ]
\]

\clearpage
\endgroup
